\documentclass[10pt,a4paper]{amsart}
\usepackage{amsmath,amssymb,amsthm,graphicx}
% Retained minimal preamble for the Stacks Project chapter more-morphisms.tex (More on Morphisms).
% Verbatim excerpt of the chapter preamble: Xy-pic setup (native preamble lines 14-18) and the native
% theorem declarations with the mathematical macros (native preamble lines 154-217). Copied mechanically;
% nothing added, removed or redefined. The wrapper class and the amsmath/amssymb/amsthm/graphicx base
% packages are supplied by the assembler, so the class/externaldocument part of the native preamble is omitted.
% For commutative diagrams we use Xy-pic
\usepackage[all]{xy}

% We use 2cell for 2-commutative diagrams.
\xyoption{2cell}
% Theorem environments.
%
\theoremstyle{plain}
\newtheorem{theorem}[subsection]{Theorem}
\newtheorem{proposition}[subsection]{Proposition}
\newtheorem{lemma}[subsection]{Lemma}

\theoremstyle{definition}
\newtheorem{definition}[subsection]{Definition}
\newtheorem{example}[subsection]{Example}
\newtheorem{exercise}[subsection]{Exercise}
\newtheorem{situation}[subsection]{Situation}

\theoremstyle{remark}
\newtheorem{remark}[subsection]{Remark}
\newtheorem{remarks}[subsection]{Remarks}

\numberwithin{equation}{subsection}

% Macros
%
\def\lim{\mathop{\mathrm{lim}}\nolimits}
\def\colim{\mathop{\mathrm{colim}}\nolimits}
\def\Spec{\mathop{\mathrm{Spec}}}
\def\Hom{\mathop{\mathrm{Hom}}\nolimits}
\def\Ext{\mathop{\mathrm{Ext}}\nolimits}
\def\SheafHom{\mathop{\mathcal{H}\!\mathit{om}}\nolimits}
\def\SheafExt{\mathop{\mathcal{E}\!\mathit{xt}}\nolimits}
\def\Sch{\mathit{Sch}}
\def\Mor{\mathop{\mathrm{Mor}}\nolimits}
\def\Ob{\mathop{\mathrm{Ob}}\nolimits}
\def\Sh{\mathop{\mathit{Sh}}\nolimits}
\def\NL{\mathop{N\!L}\nolimits}
\def\CH{\mathop{\mathrm{CH}}\nolimits}
\def\proetale{{pro\text{-}\acute{e}tale}}
\def\etale{{\acute{e}tale}}
\def\QCoh{\mathit{QCoh}}
\def\Ker{\mathop{\mathrm{Ker}}}
\def\Im{\mathop{\mathrm{Im}}}
\def\Coker{\mathop{\mathrm{Coker}}}
\def\Coim{\mathop{\mathrm{Coim}}}

% Boxtimes
%
\DeclareMathSymbol{\boxtimes}{\mathbin}{AMSa}{"02}

%
% Macros for moduli stacks/spaces
%
\def\QCohstack{\mathcal{QC}\!\mathit{oh}}
\def\Cohstack{\mathcal{C}\!\mathit{oh}}
\def\Spacesstack{\mathcal{S}\!\mathit{paces}}
\def\Quotfunctor{\mathrm{Quot}}
\def\Hilbfunctor{\mathrm{Hilb}}
\def\Curvesstack{\mathcal{C}\!\mathit{urves}}
\def\Polarizedstack{\mathcal{P}\!\mathit{olarized}}
\def\Complexesstack{\mathcal{C}\!\mathit{omplexes}}
% \Pic is the operator that assigns to X its picard group, usage \Pic(X)
% \Picardstack_{X/B} denotes the Picard stack of X over B
% \Picardfunctor_{X/B} denotes the Picard functor of X over B
\def\Pic{\mathop{\mathrm{Pic}}\nolimits}
\def\Picardstack{\mathcal{P}\!\mathit{ic}}
\def\Picardfunctor{\mathrm{Pic}}
\def\Deformationcategory{\mathcal{D}\!\mathit{ef}}
\usepackage{cleveref}
\crefname{originalreference}{original reference}{original references}
\Crefname{originalreference}{Original reference}{Original references}
\makeatletter
\newlabel{algebra-lemma-image-dense-generic-points}{{00FL}{source}{}{sourcepointer-1}{}}
\newlabel{algebra-lemma-image-dense-generic-points@cref}{{[originalreference][0][]00FL}{[1][0][]source}{}{}{}}
\makeatother
\crefname{originallemma}{original lemma}{original lemmas}
\Crefname{originallemma}{Original lemma}{Original lemmas}
\makeatletter
\newlabel{lemma-locally-principal-vertical}{{053R}{source}{}{sourcepointer-2}{}}
\newlabel{lemma-locally-principal-vertical@cref}{{[originallemma][0][]053R}{[1][0][]source}{}{}{}}
\makeatother
\crefname{originalreference}{original reference}{original references}
\Crefname{originalreference}{Original reference}{Original references}
\makeatletter
\newlabel{morphisms-lemma-dimension-fibre-at-a-point}{{02FX}{source}{}{sourcepointer-3}{}}
\newlabel{morphisms-lemma-dimension-fibre-at-a-point@cref}{{[originalreference][0][]02FX}{[1][0][]source}{}{}{}}
\makeatother
\title{A nonzerodivisor avoiding the minimal primes over the maximal ideal hits the generic point}
\author{The Stacks Project Authors}
\date{stacks-project:more-morphisms.tex:local-snapshot-20261011; GNU FDL 1.2 or later}
\begin{document}
\maketitle
\noindent\textit{Editorial note.} Counted span is the author's own proof environment, source lines 14717-14752, immediately following the statement at 14701-14715.\par\medskip
\noindent\textit{Editorial note.} The proof is by contradiction and ends with the explicit contradiction against hypothesis (3); that final step is inside the span.\par\medskip
\noindent\textit{Editorial note.} Three cross-chapter pointers (Algebra, Morphisms) plus a same-chapter lemma are recorded in assembly.external\_references.\par\medskip
\noindent\textit{Allowed standard prerequisites (editorial).}\begin{itemize}
\item Locally principal vertical parts. chapter lemma lemma-locally-principal-vertical, cited at proof lines 14744-14745. Applied to Spec(C) -> Spec(A) and Z = Spec(C/gC) to produce the contradiction that closes the proof.
\item Image of a finite map and generic points. Algebra, lemma algebra-lemma-image-dense-generic-points, cited in the proof. Controls the generic point of Spec(A) as a point of the image of the restricted morphism.
\item Dimension of the fibre at a point. Morphisms, lemma morphisms-lemma-dimension-fibre-at-a-point, cited in the proof. Supplies the fibre-dimension interpretation of the hypothesis (4) on dimensions and transcendence degrees.
\end{itemize}\medskip
\ifdefined\hypertarget\hypertarget{sourcepointer-1}{}\fi
\noindent\textit{Source reference.} Original reference 00FL: focused-seeds/geometry/S9-morph-horizontal/ref-tags.txt, source lines 4--4. Navigation pointer only: the counted proof cites this statement by \textbackslash{}ref as a standing input; the deduction the counted proof has to make is carried out inside the counted span, and no part of that deduction is delegated to this pointer.\par\medskip
\ifdefined\hypertarget\hypertarget{sourcepointer-2}{}\fi
\noindent\textit{Source reference.} Original lemma 053R: sources/stacks-more-morphisms/pinned-original-source.tex, source lines 14538--14538. Navigation pointer only: the counted proof cites this statement by \textbackslash{}ref as a standing input; the deduction the counted proof has to make is carried out inside the counted span, and no part of that deduction is delegated to this pointer.\par\medskip
\ifdefined\hypertarget\hypertarget{sourcepointer-3}{}\fi
\noindent\textit{Source reference.} Original reference 02FX: focused-seeds/geometry/S9-morph-horizontal/ref-tags.txt, source lines 5--5. Navigation pointer only: the counted proof cites this statement by \textbackslash{}ref as a standing input; the deduction the counted proof has to make is carried out inside the counted span, and no part of that deduction is delegated to this pointer.\par\medskip
\begin{lemma}
\label{lemma-horizontal}
Let $A \to B$ be a local homomorphism of local rings, and
$g \in \mathfrak m_B$. Assume
\begin{enumerate}
\item $A$ and $B$ are domains and $A \subset B$,
\item $B$ is essentially of finite type over $A$,
\item $g$ is not contained in any minimal prime over $\mathfrak m_AB$, and
\item $\dim(B/\mathfrak m_AB) +
\text{trdeg}_{\kappa(\mathfrak m_A)}(\kappa(\mathfrak m_B)) =
\text{trdeg}_A(B)$.
\end{enumerate}
Then $A \subset B/gB$, i.e., the generic point of $\Spec(A)$
is in the image of the morphism $\Spec(B/gB) \to \Spec(A)$.
\end{lemma}

\begin{proof}
Note that the two assertions are equivalent by
Algebra, Lemma \ref{algebra-lemma-image-dense-generic-points}.
To start the proof let $C$ be an $A$-algebra of finite type
and $\mathfrak q$ a prime of $C$ such that $B = C_{\mathfrak q}$.
Of course we may assume that $C$ is a domain and that $g \in C$.
After replacing $C$ by a localization we see that
$\dim(C/\mathfrak m_AC) = \dim(B/\mathfrak m_AB) +
\text{trdeg}_{\kappa(\mathfrak m_A)}(\kappa(\mathfrak m_B))$, see
Morphisms, Lemma \ref{morphisms-lemma-dimension-fibre-at-a-point}.
Setting $K$ equal to the fraction field of $A$
we see by the same reference that
$\dim(C \otimes_A K) = \text{trdeg}_A(B)$. Hence assumption
(4) means that the generic and closed fibres of the morphism
$\Spec(C) \to \Spec(A)$ have the same dimension.

\medskip\noindent
Suppose that the lemma is false. Then $(B/gB) \otimes_A K = 0$.
This means that $g \otimes 1$ is invertible in $B \otimes_A K
= C_{\mathfrak q} \otimes_A K$. As $C_{\mathfrak q}$ is a limit
of principal localizations we conclude that $g \otimes 1$
is invertible in $C_h \otimes_A K$ for some
$h \in C$, $h \not \in \mathfrak q$. Thus after replacing $C$
by $C_h$ we may assume that $(C/gC) \otimes_A K = 0$.
We do one more replacement of $C$ to make sure that the minimal
primes of $C/\mathfrak m_AC$ correspond one-to-one with the minimal
primes of $B/\mathfrak m_AB$. At this point we apply
Lemma \ref{lemma-locally-principal-vertical}
to $X = \Spec(C) \to \Spec(A) = S$ and the locally closed
subscheme $Z = \Spec(C/gC)$. Since $Z_K = \emptyset$ we see that
$Z \otimes \kappa(\mathfrak m_A)$ has to contain an irreducible
component of
$X \otimes \kappa(\mathfrak m_A) = \Spec(C/\mathfrak m_AC)$.
But this contradicts the assumption that $g$ is not contained
in any prime minimal over $\mathfrak m_AB$. The lemma follows.
\end{proof}

\end{document}
